\chapter{الگوریتم \lr{GLKM}}

\section{انگیزه و بازنویسی تابع هدف}

مقاله‌ی اصلی پروژه، \lr{GLKM}، روشی حریصانه برای بهینه‌سازی
\lr{k-means} ارائه می‌کند \cite{yuan2026}. زوجی بودن تابع جریمه
\begin{equation}
    \sum_{x\in S}\lVert x-\mu_S\rVert^2
    =
    \frac{1}{2|S|}
    \sum_{x,y\in S}\lVert x-y\rVert^2
    \label{eq:pairwise}
\end{equation}
اجازه می‌دهد هزینه‌ی هر خوشه بدون نمایش صریح مرکز آن و با استفاده از
فاصله‌های زوجی محاسبه شود. \lr{GLKM} از $n$ خوشه‌ی تک‌عضوی شروع می‌کند و در
هر مرحله زوجی را ادغام می‌کند که کمترین افزایش تابع هدف را داشته باشد؛ این
روند تا رسیدن به تعداد خوشه‌ی درخواستی $k$ ادامه می‌یابد.

\section{ماتریس فاصله}

نگهداری همه‌ی $n^2$ فاصله برای داده‌ی بزرگ مناسب نیست. مقاله یک گراف
$K$-نزدیک‌ترین همسایه می‌سازد. برای یال‌های موجود، فاصله‌ی مربع‌شده‌ی واقعی
و برای زوج‌های فاقد یال مقدار  $\gamma$ در نظر گرفته می‌شود؛
$\gamma$ بر مبنای بزرگ‌ترین فاصله‌ی موجود در گراف تنظیم می‌شود. در نتیجه،
هدف محاسبه‌شده دقیقاً تابع هدف اقلیدسی کامل نیست، بلکه تقریب آن
است. 

\section{شتاب‌دهی‌ها}

سه ایده‌ی اصلی برای کاهش هزینه‌ی جست‌وجوی ادغام‌ها به‌کار رفته است:
\begin{enumerate}
    \item به‌روزرسانی جبری هزینه‌ی ادغام با نگهداری آمار خوشه‌ها، به‌جای
          محاسبه‌ی دوباره‌ی تمام فاصله‌ها؛
    \item محدود کردن زوج‌های نامزد به همسایگی گراف:؛
    \item نگهداری نامزدهای ادغام در ساختاری مرتب برای دسترسی سریع به کمترین
          افزایش.
\end{enumerate}

\section{ایده‌ها برای توسعه این پروژه}

دو پرسش از مطالعه‌ی روش حاصل شد:
\begin{enumerate}
    \item آیا می‌توان با اجرای \lr{GLKM} روی چند مقدار $k$ و نمودار elbow
          نیاز به دریافت $k$ از کاربر را ااز بین برد؟
    \item آیا می‌توان به‌جای آغاز از $n$ خوشه‌ی تک‌عضوی، یک افراز اولیه‌ی کوچک‌تر
          ساخت تا تعداد ادغام‌ها و زمان اجرای هسته‌ی \lr{GLKM} کاهش یابد؟
\end{enumerate}
فصل‌های بعدی پیاده‌سازی و ارزیابی مقدماتی این دو پرسش را گزارش می‌کنند.
